
\begin{frame}{SGD in Hilbert Space: stochastic modified equation (SME)}
    We are interested in studying the impact on the noise on gradient dynamics, to this end they propose the following continuous time model for analysis:
    \begin{align}
        d\varphi_t = -\nabla \mathcal{F}(\varphi_t)dt + \sqrt{n} \;\sigma(\varphi_t) dB_t,
    \end{align}
    where $B_t$ is a cylindrical Wiener process (Brownian motion on $H$), and $\sigma$ is the diffusion kernel of the noise.\\ 
    Discretizing of SME (2) using Euler-Maruyama, with step-size $\eta$ (i.e., $t_k = k\eta$),
    \begin{align}
        \tilde{\varphi}_{n+1} = \tilde{\varphi}_n - \eta \nabla \mathcal{F}(\tilde{\varphi}_n) + \sqrt{\eta} \; \sigma(\tilde{\varphi}_n)( B_{(n+1)\eta} - B_{n \eta}).
    \end{align}
    \textcolor{red}{Replace S10}
\end{frame}

\begin{frame}{SGD in Hilbert Space}
    Comparing the SGD step (1) and the SME (2), we immediately see that the \textbf{drift term} matches the gradient update.\\
    And if the diffusion kernel $\sigma: U \rightarrow H$ is defined as
    \begin{align*}
        \sigma(\phi)u := \mathbb{E}_{\gamma}[V(\phi; \gamma) u(\gamma)], \quad u \in U
    \end{align*}
    where $U = L^2(\Gamma)$, the space of random variables with finite 2nd moment.\\
    $\sigma$ is a Hilbert-Schmidt operator.\\
    For this choice of $\sigma$, (3) approximates (2) with error $\mathcal{O}(\eta^2)$.\\
    % so for very small steps, the discretization of the stochastic modified equation is a good approximation of the SGD.
    \textcolor{red}{Replace S11}
\end{frame}


\begin{frame}{$\nabla \mathcal{F}$ is Local Lipschitz}
    Consider $\phi, \psi \in H$,
    \begin{align*}
        \|\nabla \mathcal{F}(\phi) - \nabla \mathcal{F}(\psi) \| &= \| \mathbb{E}[\nabla \mathcal{F}_{\gamma}(\phi) - \nabla \mathcal{F}_{\gamma}(\psi) ]\|\\
        \text{using Jensen's inequality }& \text{ and convexity of $\|\cdot\|$,}\\
        & \leq \mathbb{E}[\|\nabla \mathcal{F}_{\gamma}(\phi) - \nabla \mathcal{F}_{\gamma}(\psi)\|]\\
        \text{using assumption (4),}\qquad&\\
        & \leq \mathbb{E}[C_L(\gamma) \| \phi - \psi \|] = \mathbb{E}[C_L(\gamma)] \|\phi - \psi \|.\\
        \text{since }C_L(\gamma) \in L^2(\Gamma), \mathbb{E}[C_{L}(\gamma)^2] &< \infty,\\
        \text{using Jensen's inequality,}\quad &\\
        |\mathbb{E}[C_L(\gamma)]|^2& \leq \mathbb{E}[C_{L}(\gamma)^2] < \infty
    \end{align*}
    Hence $\mathbb{E}[C_L(\gamma)] < \infty$ and $\nabla \mathcal{F}$ is local Lipschitz with constant $\mathbb{E}[C_L(\gamma)]$.
\end{frame}



\begin{frame}{$\nabla \mathcal{F}(\phi)$ has linear growth rate}
    Compute the inner product,
    \begin{align*}
        \langle \nabla \mathcal{F}(\phi), \phi\rangle &= \langle \mathbb{E}[\nabla \mathcal{F}_{\gamma}(\phi)], \phi\rangle = \mathbb{E}[\langle \nabla \mathcal{F}_{\gamma}(\phi), \phi\rangle]
    \end{align*}\\
    \text{using assumption (4)},\\
    \begin{align*}
        \mathbb{E}[\langle \nabla \mathcal{F}_{\gamma}(\phi), \phi\rangle] &\leq \mathbb{E}[C_G(\gamma)(1 + \|\phi\|^2)] = \mathbb{E}[C_G(\gamma)](1 + \|\phi\|^2).
    \end{align*}
    Since $C_G(\gamma) \in L^1(\Gamma)$, $\mathbb{E}[C_G(\gamma)] < \infty$.\\
    Hence $\nabla \mathcal{F}(\phi)$ linear growth rate in $\|\phi\|$.
\end{frame}

\begin{frame}{$\sigma$ is Local Lipschitz}
    Let $(e_k)$ and $(f_j)$ denote orthonormal basis for the separable Hilbert spaces $U$ and $H$.\\
    For any H-S operator $T: U \rightarrow H$, Hilbert-Schmidt norm
    \begin{align*}
        \|T\|_{\mathcal{L}_2(U, H)}^2 = \sum_k \|Te_k\|^2,
    \end{align*}
    If $T = \sigma(\phi) - \sigma (\psi)$,
    \begin{align*}
        \| \sigma(\phi) - \sigma (\psi) \|_{\mathcal{L}_2(U, H)}^2 &= \sum_k \| (\sigma(\phi) - \sigma (\psi))e_k\|^2
    \end{align*}
    Recall the diffusion kernel $\sigma: U \rightarrow H$ defined as
    \begin{align*}
        \sigma(\phi)u := \mathbb{E}_{\gamma}[V(\phi; \gamma) u(\gamma)], \quad u \in U = L^2(\Gamma).
    \end{align*}
\end{frame}


\begin{frame}{$\sigma$ is Local Lipschitz}
    \begin{align*}
        \| \sigma(\phi) - \sigma (\psi) \|_{\mathcal{L}_2(U, H)}^2 =& \sum_k\|\mathbb{E}_{\gamma}[(V(\phi; \gamma) - V(\psi; \gamma)) e_k(\gamma)]\|^2\\
        \text{using Parseval's identity, } &\\
        % \text{and linearity of $\mathbb{E}[\cdot]$,}\\
        =&\sum_{j,k} \left(\mathbb{E}_{\gamma} \langle (V(\phi; \gamma) - V(\psi; \gamma))e_k(\gamma), f_j\rangle\right)^2\\
        \text{since $e_k(\gamma)$ is a scalar},\quad &\\
        =&\sum_{j,k} \left(\mathbb{E}_{\gamma} \langle (V(\phi; \gamma) - V(\psi; \gamma)), f_j\rangle e_k(\gamma) \right)^2\\
        \text{define a random variable }& a_{j}(\gamma) = \langle (V(\phi; \gamma) - V(\psi; \gamma)), f_j\rangle,\\
        \text{then } \sum_k(\mathbb{E}[a_j(\gamma) e_k])^2 &= \mathbb{E}|a_j(\gamma)|^2.\\
        =& \sum_j \mathbb{E}|\langle (V(\phi; \gamma) - V(\psi; \gamma)), f_j\rangle|^2 = \mathbb{E} \|V(\phi; \gamma) - V(\psi; \gamma)\|^2\\
        &\text{using linearity of $\mathbb{E}[\cdot]$ and parseval's identity.}
    \end{align*}
\end{frame}

\begin{frame}{$\sigma$ is Local Lipschitz}
    \begin{align*}
        \| \sigma(\phi) - \sigma (\psi) \|_{\mathcal{L}_2(U, H)}^2 = \mathbb{E} [\|V(\phi; \gamma) - V(\psi; \gamma)\|^2]
    \end{align*}
    where $V(\phi; \gamma) - V(\psi ; \gamma) = (\nabla \mathcal{F}(\phi) - \nabla \mathcal{F}(\psi)) - (\nabla \mathcal{F}_{\gamma}(\phi) - \nabla \mathcal{F}_{\gamma}(\psi)) = \mathbb{E}_{\gamma}[\nabla \mathcal{F}_{\gamma}(\phi) - \nabla \mathcal{F}_{\gamma}(\psi)] - (\nabla \mathcal{F}_{\gamma}(\phi) - \nabla \mathcal{F}_{\gamma}(\psi))$.\\
    Using the variance identity, $\mathbb{E}[\|X - \mathbb{E}[X]\|^2] = \mathbb{E}[\|X\|^2] - \|\mathbb{E}[X]\|^2 \leq \mathbb{E}[\|X\|^2]$. So,
    \begin{align*}
        \| \sigma(\phi) - \sigma (\psi) \|_{\mathcal{L}_2(U, H)}^2 &\leq \mathbb{E} [\|\nabla \mathcal{F}_{\gamma}(\phi) - \nabla \mathcal{F}_{\gamma}(\psi)\|^2]\\
        \text{using assumption (4),}\quad \quad&\\
        &\leq \mathbb{E}_{\gamma}[C_L(\gamma)^2\|\phi - \psi\|^2] = \mathbb{E}_{\gamma}[C_L(\gamma)^2] \|\phi - \psi\|^2.
    \end{align*}
    Hence, $\| \sigma(\phi) - \sigma (\psi) \|_{\mathcal{L}_2(U, H)} \leq \sqrt{\mathbb{E}_{\gamma}[C_L(\gamma)^2]} \|\phi - \psi\|$.
\end{frame}

\begin{frame}{$\sigma$ has linear growth rate}
    Using a similar approach of parseval's identity + orthonormal expansion,
    \begin{align*}
        \| \sigma(\phi) \|_{\mathcal{L}_2(U, H)}^2 &= \mathbb{E}[\|V(\phi;\gamma)\|^2] = \mathbb{E}[\| \nabla \mathcal{F}(\phi) - \nabla \mathcal{F}_{\gamma}(\psi) \|^2]\\
        \text{using the identity}& \|a - b\|^2 \leq 2 \|a\|^2 + 2\|b\|^2,\\
        &\leq 2 \|\nabla \mathcal{F}(\phi)\|^2 + 2 \mathbb{E}[\|\nabla \mathcal{F}_{\gamma}(\phi)\|^2]\\
        \text{using Jensen's: } \|\nabla &\mathcal{F}(\phi)\|^2 = \|\mathbb{E}[\nabla \mathcal{F}_{\gamma}(\phi)]\|^2 \leq \mathbb{E}[\|\nabla \mathcal{F}_{\gamma}(\phi)\|^2],\\
        \| \sigma(\phi) \|_{\mathcal{L}_2(U, H)}^2 &\leq 4\mathbb{E}[\|\nabla \mathcal{F}_{\gamma}(\phi)\|^2]
    \end{align*}
    \begin{align*}
        \text{ The norm, }\|\nabla \mathcal{F}_{\gamma}(\phi)\|^2 &= \|\nabla \mathcal{F}_{\gamma}(\phi) - \nabla \mathcal{F}_{\gamma}(0) + \nabla \mathcal{F}_{\gamma}(0)\|^2\\
        &\leq 2\|\nabla \mathcal{F}_{\gamma}(\phi) - \nabla \mathcal{F}_{\gamma}(0)\|^2 + 2\|\nabla \mathcal{F}_{\gamma}(0)\|^2\\
        &\leq 2 C_L(\gamma)^2 \|\phi\|^2 + 2\|\nabla \mathcal{F}_{\gamma}(0)\|^2
    \end{align*}
    using assumption (4).
\end{frame}
\begin{frame}{$\sigma$ has linear growth rate}
    Hence,
    \begin{align*}
        \| \sigma(\phi) \|_{\mathcal{L}_2(U, H)}^2 &\leq 4\mathbb{E}[2 C_L(\gamma)^2 \|\phi\|^2 + 2\|\nabla \mathcal{F}_{\gamma}(0)\|^2]\\
        &= 8\mathbb{E}[ C_L(\gamma)^2] \|\phi\|^2 + 8\mathbb{E}[\|\nabla \mathcal{F}_{\gamma}(0)\|^2]\\
        &\lesssim 1 + \|\phi\|^2
    \end{align*}
    i.e. there exists $c > 0$, such that $\| \sigma(\phi) \|_{\mathcal{L}_2(U, H)}^2 \leq c(1 + \|\phi\|^2)$.
\end{frame}
